\section{Introduction and the proposed analogue}
\label{s1-introduction}

The source manuscript \cite{F} constructs rational polynomials in a
variable $X$, uses local estimates to clear their coefficients, and uses
a positive integral to estimate their values at $X=\zeta(5)$. Here we
carry out the corresponding manipulations with $\zeta(9)$. The question
is what the individual formulas become, including the formulas at the
end that cease to give the desired conclusion.

We make one choice at the outset: retain the source's
$K=40n$, $N=3n$, $h=37n$, and numerator $D_N^6$. Replace the functional
and all harmonic sums, moments, and residues that depend on the zeta
order. A construction with $D_N^{10}$ would be a different experiment:
it changes the degree bookkeeping, local weights, normalization, and
real field together. It is not part of the main calculation below.

\subsection{What the final step would require}

The mechanism behind the source's main theorem can be stated without
assuming any of its estimates. Suppose integer polynomials $Q_K$ have
degree $h=\lambda K$, where $\lambda=37/40$, and satisfy
\begin{equation}
 0<Q_K(\zeta(9))\leq \exp(-\varepsilon K^2+o(K^2))
 \quad\text{for some }\varepsilon>0.
 \label{s1-target}
\end{equation}
If $\zeta(9)=a/b$ with integers $a,b$ and $b>0$, then
$b^hQ_K(a/b)$ is a positive integer. But its logarithm is at most
$\lambda K\log b-\varepsilon K^2+o(K^2)$, which tends to $-\infty$.
That is the contradiction one would seek. Positivity, integrality, and
quadratic decay have distinct jobs in this short final argument.

Our analogue writes $Q_{9,K,M}=m_{9,K,M}F_{9,K}$, with $F_{9,K}$ the
real-normalized determinant and $m_{9,K,M}>0$ the arithmetic multiplier.
For fixed $M$, the two proposed bounds are
\begin{equation}
 \log m_{9,K,M}\leq A_{9,M}K^2+o_M(K^2),\qquad
 \log F_{9,K}(\zeta(9))\leq UK^2+O(K\log K).
 \label{s1-two-bounds}
\end{equation}
The sign that matters is therefore $A_{9,M}+U$. An upper bound with a
positive coefficient does not imply that the determinant values grow;
it simply fails to prove the decay needed in \eqref{s1-target}.

\subsection{The changes to expect}

The following dictionary records the principal substitutions. The
entries in the right column are derived or qualified in the indicated
sections; the table is a guide to the calculations, not an additional
theorem.
\begin{center}
\small
\begin{tabular}{@{}>{\raggedright\arraybackslash}p{0.35\textwidth}>{\raggedright\arraybackslash}p{0.25\textwidth}>{\raggedright\arraybackslash}p{0.29\textwidth}@{}}
\toprule
Object & Source order five & Draft order nine\\
\midrule
Simple-pole coefficient of $X$ & $j^4$ & $j^8$\\
Polynomial functional & $L(P^{(3)})/4!$ & $L(P^{(7)})/8!$\\
Real weight series & $y^5\sum \ell^4e^{-2\pi\ell y}$ & $y^9\sum \ell^8e^{-2\pi\ell y}$\\
Large-prime distribution & $p^{-4}$ & $p^{-8}$\\
Ordinary-class exponent loss & $-4$ & $-8$\\
Zero-class net contribution & $+1$ & $+1$\\
First possible nonintegral moment, $p\geq11$ & $e\geq2p-3$ & $e\geq4p-5$\\
Outer unremoved-class cost & $7(\ell-2)$ & $11(\ell-2)$\\
Degree of the determinant & $h$ & $h$\\
Leading real bound $U$ & $-1.3669955$ & $-1.3669955$\\
\bottomrule
\end{tabular}
\end{center}

Two effects compete in the arithmetic. Higher harmonic sums increase
the pole denominators. Extra factors in the polynomial moments cancel
more Bernoulli denominators. Section~4 keeps both effects: in particular,
the outer polynomial correction disappears sooner, even though the
total denominator-clearing estimate becomes more costly.

The calculation in Section~5 and Appendix~B gives
\begin{equation}
 A_{9,M}=A_{5,M}+\frac{116581}{57600},\qquad
 \frac{116581}{57600}=2.0239756944\ldots .
 \label{s1-cost-shift}
\end{equation}
Here $A_{5,M}$ denotes the source's particular arithmetic bound. It is
not a claimed minimum among all normalizations. Combining the draft's
bound with $U$ leaves approximately $2.00657$ at $M=200$, and approaches
$1.97449$ as $M\to\infty$. Thus this version of the estimates has already
lost its final negative sign. Section~7 nevertheless writes out the
completion, and Appendix~C carries the perturbation and height algebra
as far as its hypotheses permit.

There is no additional restriction $9\equiv1\pmod4$ in the positive
weight calculation used here. For every odd $s\geq3$, the corresponding
weight has the positive series
\[
 \frac{2(2\pi)^{s-1}}{(s-1)!}y^s
       \sum_{\ell\geq1}\ell^{s-1}e^{-2\pi\ell y}.
\]
Both the seventh and ninth orders fit that identity. The relevant
obstruction in this draft is the numerical arithmetic cost, not a
change of sign in that series. This statement concerns the present
functional; it makes no assertion about separate modular transformation
formulas in which congruence modulo four can matter.

\subsection{Organization and limits of the exercise}

Sections~2 and 3 construct the functional and translate its local
identities. Section~4 supplies the proposed determinant exponents.
Section~5 converts them to a multiplier and a prime-sum cost.
Section~6 preserves the real energy calculation while changing the
weight and its subquadratic errors. Section~7 displays the failed sign
comparison. Appendix~A shows how the source's comparison measure is
reused; Appendix~B gives the new rational arithmetic totals;
Appendix~C develops perturbation, rational-approximation, and height
formulas, retaining a hypothetical decay parameter only where needed.

The inner completion step remains marked as a formal transplant. The
source's real interval certificate is reused, not reverified. These
qualifications do not prevent one from examining the substitutions,
but they do prevent treating the draft as a checked theorem. The
positive combined upper bound is also specific to this construction
and these estimates; it rules out neither improved estimates nor a
different construction.
